\section{Pretraining contamination analysis}
\label{app:contamination}

This appendix is load-bearing for the paper's central claim. It states
the contamination question precisely, lays out the available evidence,
and identifies what EXP-007 / EXP-008 can and cannot conclude from that
evidence.

\subsection{The question}
The EXP-007 result --- frozen DermLIP plus a linear scaffold reaches
test $\auroc = 0.944 \pm 0.003$ on \strongholdtwo\ --- is consistent with
two distinct hypotheses:
\begin{enumerate}[label=(\roman*), leftmargin=2em]
\item \emph{Dermoscopy-domain transfer.} Pretraining on dermatology image-
      text pairs (Derm1M) organises the embedding space such that a
      single linear direction captures stable-pair similarity uniformly
      across the three nuisance families, including the unseen
      \strongholdtwo.
\item \emph{HAM10000 image-level overlap.} Derm1M contains HAM10000 (or
      near-duplicate) images via its uncontrolled literature/web sources,
      so DermLIP's embedding space is well-fit to HAM10000 lesion identity
      in a way that survives our synthetic post-hoc nuisance.
\end{enumerate}
Hypotheses (i) and (ii) are not mutually exclusive: a positive
EXP-009 outcome (\S\ref{sec:followup}) would support (i) being load-bearing;
a low EXP-009 outcome would push the explanation toward (ii).

\subsection{Is HAM10000 in Derm1M?}
Derm1M \citep{yan2025derm1m} names exactly two public datasets as image
sources --- SCIN and MSKCC (an ISIC-archive subset \citep{codella2018isic})
--- and HAM10000 \citep{tschandl2018ham10000} is not among them: in the
Derm1M paper HAM10000 appears only as a downstream evaluation benchmark,
as a comparison row in the related-datasets table, and as a source of
disease \emph{label names} for the standardised taxonomy. At the level of
named sources, HAM10000 is therefore not a constituent of Derm1M, and the
named MSKCC source is a Memorial-Sloan-Kettering ISIC subset contributed by
a different institution than HAM10000's Vienna/Rosendahl images, so it is
unlikely to introduce HAM10000 images directly.

Overlap is nonetheless present through Derm1M's \emph{uncontrolled} sources.
Derm1M documents no deduplication or decontamination against HAM10000 or any
downstream benchmark, and beyond its two named public datasets it
bulk-aggregates PubMed figures, educational video, and social-media
dermatology content, in which HAM10000 --- the most widely reproduced public
dermoscopy dataset --- circulates freely. A perceptual-hash audit of HAM10000
against the (access-gated) Derm1M corpus confirms this: at least
\textbf{13 distinct HAM10000 images}, drawn from at least four PubMed Central
articles, are present in Derm1M's PubMed/literature partition, verified by
agreement across four independent perceptual hashes (pHash, dHash, aHash,
wHash) and by visual inspection (one journal figure alone, PMC9777332, is a
multi-panel montage reproducing roughly five HAM10000 dermoscopy images). The
13 is a lower bound: dermoscopy images collide heavily under any single
perceptual hash, so only multi-hash-plus-visual matches were counted; the
forum and textbook candidates were not all individually verified; and the
$16.6$\,GB YouTube-frame partition was not audited. DermLIP's pretraining was
therefore \emph{not} HAM10000-free. The confirmed scale is small relative to
HAM10000's $10{,}015$ images, which makes wholesale memorisation an unlikely
\emph{sole} explanation of the EXP-007 win, but it does not clear the
confound; the EXP-009 partition (\S\ref{sec:followup}), which holds
``HAM10000 absent'' constant by construction, remains the way to separate
dermoscopy-domain transfer from overlap. Full method and per-article findings
are in the repository's overlap-audit report.

\subsection{What EXP-007 alone can conclude}
EXP-007 demonstrates a $+0.66$ test $\auroc$ swing relative to
EXP-006b (OpenAI CLIP) when the only varied axis is the pretraining
data. EXP-007 alone cannot distinguish (i) from (ii), and the paper
explicitly does not attribute the swing to dermoscopy-domain transfer
on the basis of EXP-007 alone.

\subsection{What EXP-008 adds and what it does not}
EXP-008 introduces a third pretraining-data point: BiomedCLIP's PMC-15M,
which is approximately $15\times$ larger than Derm1M but is composed of
biomedical figure--caption pairs with a very low fraction of
dermoscopy and no documented raw HAM10000/ISIC archive ingestion. The result --- test
$\auroc = 0.329 \pm 0.012$ --- contributes three readings:
\begin{itemize}
\item \textbf{Scale alone is not the explanation.} Fifteen million
      biomedical pairs without documented raw HAM10000/ISIC archive
      ingestion do not unlock the proxy, so the
      EXP-007 win is not attributable to ``more pretraining data,''
      ``more medical pretraining data,'' or any monotonic function of
      pretraining-corpus size.
\item \textbf{General medical pretraining is not the explanation.} The
      $+0.04$ web $\to$ general-medical step is small relative to the
      $+0.62$ general-medical $\to$ dermoscopy step. The unlock is
      localised at the dermoscopy-specific corpus boundary.
\item \textbf{Dermoscopy-vs-HAM10000 remains unresolved.} EXP-008's
      pretraining corpus differs from EXP-007's in two ways
      simultaneously --- domain narrowness and (likely) HAM10000
      overlap (confirmed, small-scale) --- so
      EXP-008 cannot separate the two factors. EXP-009
      (\S\ref{sec:followup}) is designed precisely to break this
      confound by holding ``HAM10000 absent'' constant while varying
      the pretraining-domain narrowness from PMC-15M to dermoscopy.
\end{itemize}

\subsection{Boundary of inference}
The strongest claim the paper makes about the EXP-007 swing is the
following: \emph{under the protocol of \S\ref{sec:method}, varying only
the pretraining-data domain among DermLIP, BiomedCLIP, and OpenAI
CLIP --- holding architecture, predictor, optimiser, hyperparameters,
and data pipeline constant --- produces a $+0.66$ $\auroc$ swing
localised to the dermoscopy-specific corpus.} The claim is a
characterisation of the empirical gradient, not an attribution to
domain transfer. The corresponding negative claim --- ``DermLIP
generalises to \strongholdtwo\ because dermoscopy pretraining unlocks
nuisance invariance'' --- is not made anywhere in the paper.
